\section{Dynamics: payoffs, replicator dynamics, rest points and ESS}\label{ch:dynamics}

\subsection{Payoffs}\label{sec:payoffs}\emph{module} \code{engine/payoffs.py}

\subsubsection{\texttt{Payoffs}}

\whatit
Turns score distributions into expected payoffs: the probability that a Can of one strategy helps a Need of another, and
from it each strategy's payoff in the donation game.

\howit
Let $P_k(c)$ be the fraction of strategy $k$ whose score is in class $c$. The probability that a Can of strategy $i$
helps a Need of strategy $k$ is
\[
 H_{ik}=\sum_{c}\sum_{d} P_k(c)\;h_i(c,d)\;w^{(i)}_d ,
\]
computed with one \code{einsum} over classes and scores. A Can of $i$ helping a Need of $k$ costs $c$ to $i$ and gives $b$
to $k$, so with the role factor $\kappa$
\[
 \pi_i=\kappa\Bigl(b\sum_k\rho_k H_{ki}-c\sum_k\rho_k H_{ik}\Bigr)
\]
(``helped by Cans of every strategy'' minus ``cost of helping Needs of every strategy''); this is the payoff per
interaction, role factor included.

\paramsh
\begin{params}
benefit & float, \code{2.0} & $b$, what a receiver gains.\\
cost & float, \code{1.0} & $c$, what a donor pays.\\
role\_factor & float, \code{0.25} & $\kappa$.\\
population, w & & \code{help\_probabilities(population, w)} returns $H$; \code{payoffs(population, w)} returns $\pi$ ($w$: the $n\times L$ score distributions).\\
\end{params}

\exh
\ex{payoffs}

%--------------------------------------------------------------------------
\subsection{Replicator dynamics}\label{sec:replicator}\emph{module} \code{engine/replicator.py}

\subsubsection{\texttt{ReplicatorDynamics} and \texttt{Trajectory}}

\whatit
Evolves the strategy frequencies by the replicator equation
\[
 \dot\rho_i=\rho_i\,(\pi_i-\bar\pi),\qquad \pi=A\rho,\ \ \bar\pi=\rho\cdot\pi,\qquad A_{ik}=\kappa\,(b\,H_{ki}-c\,H_{ik}),
\]
with the score distributions evolving underneath, and returns the path as a \code{Trajectory}.

\howit
\code{payoff\_matrix(population, w)} gives $A$ and \code{velocity(population, w)} gives $(\dot\rho,\pi)$. \code{run} then
repeats, for \code{n\_steps} $+1$ states: solve the score chains for the current $\rho$ (section~\ref{sec:solver}),
record $\rho$, the payoffs, the fraction at score 0 and the mean score, compute $\dot\rho$, and take an Euler step
$\rho\leftarrow\rho+\mathrm{dt}\,\dot\rho$, clipped at $0$ (a strategy that overshoots is extinct) and renormalised.
It stops early at a \emph{rest point}, when $\max|\dot\rho|\le$ \code{rest\_tol}. Time is \code{step}$\cdot$\code{dt}.
At every step the chains are restarted from \code{w\_init}, so the payoffs are a function of $\rho$ alone; when the chains
do not converge (\code{scores\_converged} is false) the payoffs are those after the solver's \code{max\_iter} steps and so
depend on \code{w\_init} and \code{max\_iter}. With \code{warm\_start=True} each step continues from the previous step's
distributions instead.

\paramsh
\begin{params}
solver & \code{StableScoreSolver} & solves the score chains at each step; its \code{tol} and \code{max\_iter} apply.\\
payoffs & \code{Payoffs} & benefit, cost and role factor.\\
dt & float, \code{0.1} & the Euler step of $\rho$.\\
warm\_start & bool, \code{False} & continue the score chains from the previous step.\\
rest\_tol & float, \code{1e-9} & stop when $\max|\dot\rho|$ falls to this.\\
strategies, rho, w\_init & & \code{run(strategies, rho, w\_init, n\_steps, record\_scores=False, progress=None)}: the strategies, the starting frequencies, the initial score distribution(s), the number of steps; \code{record\_scores} keeps every score distribution in \code{Trajectory.w}; \code{progress(done, total)} is called after each step.\\
\end{params}
\code{Trajectory}: \code{t}, \code{rho}, \code{payoffs}, \code{w0}, \code{mean\_score} (one row per recorded state),
\code{scores\_converged}, \code{at\_rest}, \code{w}, \code{mean\_payoff}.

\exh
\ex{replicator}

%--------------------------------------------------------------------------
\subsection{The strategy simplex and its rest points}\label{sec:simplex}\emph{module} \code{engine/simplex.py}

\subsubsection{\texttt{ReplicatorField}, \texttt{PointResult}}

\whatit
Evaluates the replicator's velocity $\dot\rho$ at \emph{any} point $\rho$ of the simplex, with the payoffs, helping
probabilities and score distributions behind it, one point at a time or many in parallel. This is what the arrows of
the phase portrait are made of.

\howit
\code{evaluate(rho)} solves the score chains for that $\rho$ (restarting from \code{w\_init}, as \code{run} does), then
computes $A$, $\pi$, $\dot\rho=\rho(\pi-\rho\cdot\pi)$ and $H$ and returns them in a \code{PointResult}.
\code{evaluate\_many} spreads the points over worker processes (the field is sent to each worker once, so it must be picklable;
with the \code{spawn} start method the calling script needs an \code{if \_\_name\_\_ == "\_\_main\_\_":} guard).

\paramsh
\begin{params}
dynamics & \code{ReplicatorDynamics} & provides the solver and the payoffs.\\
strategies & list of \code{Strategy} & the $n$ strategies.\\
w\_init & array-like & the initial score distribution(s).\\
rho & array of $n$ floats & \code{evaluate(rho)}, \code{velocity(rho)}.\\
points, n\_jobs, progress & list, int or \code{None}, callable & \code{evaluate\_many}: the points; the number of processes (\code{None}: one per CPU, 1: this process); \code{progress(done, total)}.\\
\end{params}
\code{PointResult}: \code{rho}, \code{velocity}, \code{payoffs}, \code{payoff\_matrix}, \code{help\_matrix}, \code{w}, \code{converged}.
\code{velocity(rho)} returns \code{(velocity, converged)}.

\subsubsection{\texttt{simplex\_grid}}

\whatit
A regular grid of points on the simplex.

\howit
All $n$ non-negative integer counts that sum to \code{resolution}, divided by \code{resolution}. For $n=2$ the points run evenly
from $e_1$ to $e_2$.

\paramsh
\begin{params}
n & int $\ge1$ & the number of strategies.\\
resolution & int $\ge1$ & points per side.\\
\end{params}

\subsubsection{\texttt{find\_rest\_points} and \texttt{RestPoint}}

\whatit
Finds the rest points of a velocity field on the simplex (for $n=2$ or $3$) and says whether each is stable, unstable, a
saddle or neutral.

\howit
Vertices are always rest points. On an edge, a rest point is a sign change of $\pi_b-\pi_a=v_b/(\rho_a\rho_b)$ between
grid points, refined by Brent's method. Inside the triangle ($n=3$) a rest point starts from a grid cell where
$\pi_1-\pi_3$ and $\pi_2-\pi_3$ both change sign and is refined by Newton's method. Stability comes from the
eigenvalues of the Jacobian of the replicator restricted to the simplex, taken by finite differences placed so that
every point evaluated stays in the simplex: along a face, central differences; off the face, forward differences, which
equal the \emph{invasion fitness} $\pi_k-\bar\pi$ of the missing strategy $k$. The type is \code{stable} if every
real part is below $-$tolerance, \code{unstable} if every one is above it, \code{saddle} if both occur and
\code{neutral} otherwise. Segments or curves of rest points are not reported, except where a curve meets an edge.

\paramsh
\begin{params}
velocity\_fn & callable & \code{velocity\_fn(rho)} $\to$ \code{(velocity, converged)}, e.g. \code{field.velocity}.\\
n & 2 or 3 & the number of strategies.\\
resolution & int, \code{20} & grid used if \code{grid} is not given.\\
grid, grid\_velocities, grid\_converged & arrays, optional & a grid already evaluated (any resolution, any order), to avoid evaluating it again.\\
xtol & float, \code{1e-6} & precision of the root finding.\\
neutral\_tol & float, \code{1e-7} & eigenvalues smaller than this (or than twice their estimated error) count as zero; use about ten times the solver's \code{tol} for engine fields.\\
\end{params}
Each \code{RestPoint} has \code{rho}, \code{kind} (\code{stable}, \code{unstable}, \code{saddle}, \code{neutral}),
\code{eigenvalues}, \code{location} (\code{vertex}, \code{edge}, \code{interior}) and \code{converged}.

\exh
\ex{field}

%--------------------------------------------------------------------------
\subsection{Evolutionarily stable strategies}\label{sec:ess}\emph{module} \code{engine/ess.py}

\subsubsection{\texttt{ESSFinder}}

\whatit
Picks out, from a list of candidate strategies, those that are evolutionarily stable against a list of invaders.

\howit
Maynard Smith and Price: $A$ is an ESS if for every other strategy $B$ either $E(A,A)>E(B,A)$ (a rare $B$ does worse
than the common $A$), or $E(A,A)=E(B,A)$ and $E(A,B)>E(B,B)$. The finder tests this with a resident $A$ at frequency
$\rho$ (default $0.99$) against an invader $B$ at $1-\rho$, using the payoffs of the two-strategy population
(section~\ref{sec:payoffs}, from the stable score distributions at the final step):
\begin{enumerate}
\item Skip a candidate that earns almost nothing alone ($4\pi<$ \code{fitness\_threshold}): never-helping populations
are not interesting ESS.
\item For every invader $B$ of another name: if $A$'s payoff is higher, go on; if lower, $A$ is not an ESS; if equal,
swap the frequencies and require that $A$ does not do worse than $B$.
\item Finally require $\pi_A\ge10^{-5}$ for the candidate alone with the solver restricted to \code{final\_max\_iter} steps.
\end{enumerate}

\paramsh
\begin{params}
solver & \code{StableScoreSolver} & used for every population.\\
payoffs & \code{Payoffs} & benefit, cost, role factor.\\
rho & float, \code{0.99} & frequency of the resident when tested against an invader.\\
fitness\_threshold & float, \code{1e-5} & candidates whose $4\times$ payoff alone is below this are skipped.\\
final\_max\_iter & int, \code{100} & \code{max\_iter} of the last check.\\
candidates, invaders, w\_init & & \code{finder(candidates, invaders, w\_init)} returns the candidates that are ESS against every invader; \code{is\_ess(A, invaders, w\_init)} tests one.\\
\end{params}
Also \code{payoff\_alone(strategy, w\_init)} and \code{payoffs\_against(A, B, rho\_A, w\_init)}.

\exh
\ex{ess}
